\scanpage{11}\mspage{9}
Comme les foncteurs canoniques
\diagram{\begin{tikzcd}[ampersand replacement=\&,column sep=small]
\mathrm{Ind}\,\mathrm{Cat}\,\mathrm{hot} \arrow[r,"{\mathrm{Ind}\,\unsure{\varphi|\psi|\phi}}"] \& \mathrm{Ind}\,\mathrm{Hot} \arrow[r,"{\mathrm{Ind}\,\unsure{\tau|+|\gamma}}"] \& \mathrm{Ind}\,\mathrm{Hot}_0^{\wedge} \arrow[r,"{\varinjlim}"] \& \mathrm{Hot}_0^{\wedge}\\
\mathrm{Cat}\,\mathrm{hot} \arrow[u,hook] \arrow[r,"{\varphi}"] \& \mathrm{Hot} \arrow[u,hook] \arrow[r,"{\tau}"] \& \mathrm{Hot}_0^{\wedge} \arrow[u,hook] \arrow[ur,"{\mathrm{id}}"'] \&
\end{tikzcd}}
commutent aux produits finis, toute théorie \unsure{de|des} ``Hom externes'' à valeurs dans une de ces catégories, en définit une à valeurs dans chacune de celles dans lesquelles elle s'envoie.
\deleted{\unsure{Aussi}} D'ailleurs, ces foncteurs \deleted{sont} ne sont pas tous pl. fid. (c'est en tous cas faux pour $\varphi$ et $\mathrm{Ind}\,\varphi$, et j'ignore ce qu'il en est de $\tau$ et $\mathrm{Ind}\,\tau$ -- les flèches verticales par contre sont pl. fid.)
% slanted
\margin{et \unsure{églmt|égt} $\mathrm{Ind}\,\mathrm{Hot}_0^{\wedge}\xrightarrow{\varinjlim}\mathrm{Hot}_0^{\wedge}$ n'est \uline{pas} pl. fidèle}
Mais ils induisent des bijections pour les $\pi_0$. Donc la catégorie $M_{0\pi}$ déduite d'une théorie des Hom externes \inserted{à val.} dans une de ces catégories, et celle déduite dans une autre où elle s'envoie, est la $=$ à isom. can. près. 
Pour ce qui est de cette définition \inserted{de $M_{0\pi}$}, rien n'est changé donc quand on regarde les $\mathrm{CH}_{\infty}$ comme des objets de $\mathrm{Ind}\,\mathrm{Hot}$, ou seulement de $\mathrm{Ind}\,\mathrm{Hot}_0^{\wedge}$\inserted{ou $=$ \deleted{\unsure{slt}} de $\mathrm{Hot}_0^{\wedge}$}, ou quand on regarde les $\mathrm{Ch}_{\infty}$ comme des objets
% archive sheet number: 10

